% =========================================================
%  sesi/05/sesi.tex
%  Sesi 5 — Evaluasi, Playtest, dan Presentasi
% =========================================================

\sesi{05}{Evaluasi, Playtest, dan Presentasi}{4.3}{150 menit}

Ini pertemuan minggu kelima belas, dan pertemuan praktikum yang terakhir.
Level Rak sudah lengkap: mekanik, animasi, antarmuka, dan bunyi telah
tersambung. Yang belum dikerjakan adalah memeriksa apakah seluruh bagian itu
benar-benar berjalan ketika orang lain memainkannya.

Sesi ini tidak menambah fitur. Pekerjaannya adalah mengubah cara Anda memandang
karya sendiri: dari pembuat menjadi penguji. Anda menyusun
\istilah{checklist} fungsional dari \istilah{GDD}, membuat \istilah{build} yang
dapat dijalankan di komputer lain, melaksanakan \istilah{playtest} antar
kelompok, mencatat temuan pada \istilah{log bug}, memeriksa kesesuaian GDD
dengan build, lalu menyusun naskah demo lima sampai sepuluh menit.

Seluruh dokumen yang dihasilkan pada sesi ini menjadi bagian dari Proyek Akhir.
Pada minggu keenam belas, UAS menilai proyek yang dihasilkan modul ini, yaitu
build beserta dokumen pendukungnya.

% ---------------------------------------------------------
\bagian{Tujuan Sesi}

\begin{tujuanbox}
Setelah menuntaskan sesi ini, Anda dapat menunjukkan:
\begin{enumerate}[leftmargin=1.5em, itemsep=0.2em]
  \item \istilah{build} \NamaPermainan\ yang berjalan pada komputer yang tidak
        memiliki Unity;
  \item \istilah{GDD} satu halaman yang memuat inti permainan, kontrol, aturan
        menang dan kalah, serta batas lingkup;
  \item \istilah{checklist} fungsional 20--30 butir yang diturunkan dari GDD
        dan terbagi menjadi mekanik, integrasi, dan keseluruhan;
  \item lembar \istilah{playtest} yang terisi oleh sekurang-kurangnya dua
        penguji dari kelompok lain, beserta waktu dan titik berhentinya;
  \item \istilah{log bug} yang setiap catatannya dapat diulang dan memiliki
        berat serta status penanganan;
  \item naskah demo lima sampai sepuluh menit beserta refleksi perbandingan
        pekerjaan Python dan Unity.
\end{enumerate}
\end{tujuanbox}

% ---------------------------------------------------------
\bagian{Prasyarat}

\begin{prasyaratbox}
\begin{itemize}[leftmargin=1.4em, itemsep=0.25em]
  \item Checkpoint sesi 4 sudah berjalan: pemain beranimasi, skor dan nyawa
        tampil, bunyi terdengar pada peristiwanya, dan pintu dapat dibuka.
  \item Anda telah membaca Bab XVI bagian 16.1 sampai 16.6, yaitu bagian
        tentang pengujian, \istilah{checklist} fungsional, pemeriksaan
        otomatis, lembar \istilah{playtest}, \istilah{log bug}, dan
        presentasi.
  \item GDD kelompok Anda dari awal perkuliahan sudah tersedia, walaupun masih
        berupa rancangan.
  \item Dua kelompok lain bersedia menjadi penguji, dan satu komputer lain
        tanpa Unity tersedia untuk menguji build.
\end{itemize}
\end{prasyaratbox}

% ---------------------------------------------------------
\bagian{Peta Kemampuan}

\petasesi
  {Level lengkap yang hanya pernah diuji oleh pembuatnya sendiri.}
  {\istilah{Checklist} fungsional, \istilah{playtest} antar kelompok,
   \istilah{log bug}, kesesuaian GDD dengan \istilah{build}, dan naskah demo.}
  {UAS minggu 16: uji fungsional dosen atas build dan dokumen modul ini.}

% ---------------------------------------------------------
\bagian{Konsep yang Dipakai}

\subsection*{Pembuat tidak dapat menilai karyanya sendiri}

Anda sudah tahu tombol mana yang harus ditekan dan jalur mana yang harus
dilewati. Pengetahuan itu tidak dimiliki pemain baru, dan justru pengetahuan
itulah yang menutupi masalah. Level yang tampak jelas bagi pembuatnya dapat
membingungkan pemain lain selama beberapa menit pertama. Karena itu, pengujian
menuntut orang lain, dan menuntut Anda menahan diri untuk tidak menolong
mereka.

\subsection*{Checklist adalah janji, bukan keinginan}

\istilah{Checklist} fungsional berisi pernyataan yang harus bernilai benar
sebelum build dianggap selesai. Setiap butirnya diturunkan dari GDD, lalu
diubah menjadi pemeriksaan yang dapat dijawab lulus atau tidak lulus. Butirnya
dikelompokkan menjadi tiga: mekanik, integrasi, dan keseluruhan. Kelompok
ketiga hanya dapat diperiksa pada build yang dijalankan dari awal sampai akhir.

\subsection*{Playtest mengamati, bukan bertanya}

Pertanyaan ``apakah permainannya bagus'' menghasilkan jawaban sopan yang tidak
dapat ditindaklanjuti. Yang berguna adalah catatan atas apa yang benar-benar
dilakukan penguji: di mana ia berhenti, berapa lama, dan tindakan apa yang
membuatnya bingung. Catat waktu dan tempat, bukan hanya kesan.

\subsection*{Log bug harus dapat diulang}

Catatan yang berguna memuat tiga hal: apa yang terjadi, apa yang seharusnya
terjadi, dan urutan tindakan yang menyebabkannya. ``Bunyinya kadang tidak
terdengar'' tidak dapat diperbaiki; ``bunyi artefak ketiga tidak terdengar
setelah dua artefak lain diambil'' dapat diperbaiki pada hari yang sama.
Setiap catatan ditimbang beratnya dan diberi status.

\subsection*{GDD dan build boleh berbeda}

Perbedaan antara GDD dan build hampir pasti ada. Yang salah adalah perbedaan
yang tidak disadari. Tiga keadaannya: sesuai, GDD diperbarui karena perubahan
yang disengaja, dan belum dikerjakan. Janji yang belum dikerjakan dicantumkan
sebagai kekurangan, bukan dihapus dari GDD agar laporannya tampak lengkap.

% ---------------------------------------------------------
\bagian{Tutorial}

Nama menu ditulis seperti yang tampak pada Unity LTS. Kerjakan berurutan;
jangan melompat ke presentasi sebelum build dan lembar pengujian selesai.

\begin{langkah}
  \item \textbf{Bekukan fitur dan buat build.} Simpan scene, lalu pilih
        \menu{File > Build Settings}. Tekan \menu{Add Open Scenes} agar
        \file{LevelRak} masuk ke daftar. Pilih platform \menu{Windows},
        tentukan folder keluaran, lalu tekan \menu{Build}. Jangan memakai
        \menu{Build And Run}; Anda perlu berkas hasilnya, bukan menjalankannya
        di komputer sendiri.

  \item \textbf{Uji build di komputer lain.} Salin seluruh isi folder keluaran,
        bukan hanya berkas \file{.exe}. Jalankan pada komputer yang tidak
        memiliki Unity. Periksa layar, suara, dan kontrol. Build yang hanya
        berjalan di komputer pembuatnya belum memenuhi syarat.

  \item \textbf{Tulis GDD satu halaman.} Buat berkas \file{gdd.md} yang memuat
        identitas permainan, inti permainan dalam satu paragraf, tabel kontrol,
        aturan menang dan kalah, isi level, umpan balik yang terlihat dan
        terdengar, serta hal yang sengaja di luar lingkup.

  \item \textbf{Turunkan checklist.} Baca kembali GDD, lalu tulis 20--30 butir
        yang dapat diperiksa. Bagi menjadi tiga kelompok: mekanik, integrasi,
        dan keseluruhan. Simpan sebagai \file{checklist-fungsional.md} dengan
        kolom putaran 1 dan putaran 2.

  \item \textbf{Jalankan putaran pertama.} Mainkan build dari awal sampai
        pintu, lalu isi setiap butir dengan lulus atau tidak lulus. Jangan
        memperbaiki apa pun selama putaran ini; catat saja.

  \item \textbf{Susun lembar playtest.} Buat \file{lembar-playtest.md} yang
        memuat keterangan penguji, build yang diuji, waktu penyelesaian, titik
        berhenti, kesalahan yang terlihat, masukan penguji, dan catatan
        pengamat. Tuliskan satu tujuan sesi: bagian mana yang paling Anda
        ragukan.

  \item \textbf{Laksanakan playtest.} Serahkan kendali kepada dua penguji dari
        kelompok lain tanpa memberi petunjuk. Perkenalkan hanya tujuannya,
        lalu diam. Catat waktu dan tempat setiap kali penguji berhenti lebih
        dari sepuluh detik. Jangan membela karya Anda selama sesi berlangsung.

  \item \textbf{Pindahkan temuan ke log bug.} Untuk setiap temuan yang dapat
        diulang, tulis urutan tindakan, hasil yang salah, hasil yang
        seharusnya, berat, dan status. Simpan sebagai \file{log-bug.md}.
        Temuan yang belum dapat diulang disempitkan dahulu, bukan dicatat
        sebagai ``kadang-kadang''.

  \item \textbf{Perbaiki temuan berat satu per satu.} Setelah setiap
        perbaikan, jalankan kembali checklist fungsional dan isi kolom putaran
        2. Menggabungkan beberapa perbaikan membuat Anda tidak tahu perbaikan
        mana yang bekerja.

  \item \textbf{Bandingkan GDD dengan build.} Tandai setiap janji GDD sebagai
        sesuai, diperbarui, atau belum dikerjakan. Perbarui GDD untuk
        perubahan yang disengaja, dan cantumkan yang belum dikerjakan sebagai
        kekurangan yang diketahui.

  \item \textbf{Tulis naskah demo.} Susun \file{naskah-demo.md} untuk lima
        sampai sepuluh menit: pembukaan satu menit, demo langsung empat sampai
        enam menit, temuan pengujian satu sampai dua menit, dan penutup satu
        menit. Latih dengan pengukur waktu, lalu siapkan video dan build
        cadangan.

  \item \textbf{Sajikan dan refleksikan.} Tampilkan build, bukan scene di
        editor. Sesudah presentasi, tulis satu paragraf refleksi perbandingan
        pekerjaan yang Anda tulis sendiri pada Python dan pekerjaan yang
        diserahkan kepada Unity.
\end{langkah}

\begin{ujibox}
Setelah Langkah 12, periksa lima hal berikut.
\begin{itemize}[leftmargin=1.4em, itemsep=0.2em]
  \item Build berjalan pada komputer lain tanpa Unity, lengkap dengan suara.
  \item Checklist memuat 20--30 butir dari ketiga kelompok, dan kedua kolom
        putaran terisi.
  \item Lembar playtest memuat dua penguji, waktu penyelesaian, dan titik
        berhenti yang dapat disebutkan kembali.
  \item Setiap catatan log bug memuat urutan tindakan yang dapat diulang,
        berat, dan status.
  \item Naskah demo dapat diselesaikan di bawah sembilan menit.
\end{itemize}
\end{ujibox}

% ---------------------------------------------------------
\bagian{Kode dan Dokumen Sesi Ini}

Sesi ini tidak menambah skrip C\# baru. Seluruh skrip dari sesi 4 disalin ke
folder checkpoint, dan yang berubah hanyalah nilai pada \menu{Inspector} serta
beberapa baris perbaikan yang tercatat pada \istilah{log bug}. Keluaran sesi
ini adalah dokumen, bukan kode.

{\footnotesize
\renewcommand{\arraystretch}{1.25}
\begin{tabularx}{\linewidth}{@{}L{4.3cm} Y@{}}
\toprule
\textbf{Berkas} & \textbf{Isi} \\
\midrule
\file{Scripts/} & Tujuh skrip sesi 4: kontrol, animasi, artefak, pintu, HUD,
  pemusat bunyi, dan pemicu bunyi \\
\file{scene.md} & Susunan scene dan nilai Inspector yang menjadi dasar build \\
\file{gdd.md} & GDD satu halaman: inti, kontrol, aturan, lingkup \\
\file{checklist-fungsional.md} & 20--30 butir dari GDD beserta hasil dua
  putaran \\
\file{lembar-playtest.md} & Catatan dua penguji: waktu, titik berhenti,
  kesalahan, masukan \\
\file{log-bug.md} & Temuan yang dapat diulang, berat, dan status \\
\file{naskah-demo.md} & Susunan demo lima sampai sepuluh menit \\
\bottomrule
\end{tabularx}
}

\subsection*{Contoh GDD satu halaman}

Berkas berikut adalah bentuk GDD yang dipakai pada checkpoint. Perhatikan bahwa
isinya muat pada satu halaman dan hanya memuat hal yang dapat diperiksa pada
build.

\tampilkode[basicstyle=\ttfamily\scriptsize\color{ubCodeText}]
  {code/sesi-05/gdd.md}
  {gdd.md: GDD satu halaman yang menjadi dasar checklist}

\subsection*{Tabel checklist fungsional}

Tabel berikut adalah bentuk yang dipakai pada checkpoint. Setiap butir
diturunkan dari satu pernyataan GDD, dan hasilnya diisi pada setiap putaran.

{\footnotesize
\renewcommand{\arraystretch}{1.2}
\begin{tabularx}{\linewidth}{@{}C{0.6cm} L{1.9cm} Y C{0.8cm} C{0.8cm}@{}}
\toprule
\textbf{No} & \textbf{Kelompok} & \textbf{Butir pemeriksaan} &
\textbf{P1} & \textbf{P2} \\
\midrule
1 & Mekanik & Pemain berjalan ke kiri dan kanan & L & L \\
2 & Mekanik & Pemain melompat hanya ketika menyentuh tanah & L & L \\
3 & Mekanik & Pemain tidak dapat melompat berkali-kali di udara & T & L \\
4 & Mekanik & Artefak lenyap dan skor bertambah ketika disentuh & L & L \\
5 & Integrasi & Animasi berganti menurut keadaan pemain & L & T \\
6 & Integrasi & Bunyi ambil terdengar pada setiap artefak & T & L \\
7 & Keseluruhan & Level dapat diselesaikan dari awal sampai pintu & L & L \\
8 & Keseluruhan & Satu penyelesaian berlangsung 5--10 menit & L & L \\
\bottomrule
\end{tabularx}
}

\subsection*{Tabel lembar playtest}

{\footnotesize
\renewcommand{\arraystretch}{1.2}
\begin{tabularx}{\linewidth}{@{}L{4.0cm} Y@{}}
\toprule
\textbf{Bagian} & \textbf{Isi yang dicatat} \\
\midrule
Keterangan penguji & Nama atau kelompok, dan pengalaman bermain sejenis \\
Berkas yang diuji & Nama build dan tanggalnya \\
Waktu penyelesaian & Lama permainan, atau keterangan berhenti sebelum selesai \\
Titik berhenti & Tempat dan lama berhenti lebih dari sepuluh detik \\
Kesalahan yang terlihat & Keadaan yang tidak seharusnya terjadi, beserta
  urutan tindakannya \\
Masukan penguji & Pendapat penguji, dicatat apa adanya \\
Catatan pengamat & Hal yang Anda lihat tetapi tidak disebutkan penguji \\
\bottomrule
\end{tabularx}
}

\subsection*{Tabel log bug}

{\footnotesize
\renewcommand{\arraystretch}{1.2}
\begin{tabularx}{\linewidth}{@{}C{0.6cm} L{1.8cm} Y C{1.1cm} C{1.3cm}@{}}
\toprule
\textbf{No} & \textbf{Jenis} & \textbf{Urutan yang menyebabkannya dan hasil
yang salah} & \textbf{Berat} & \textbf{Status} \\
\midrule
1 & Integrasi & Ambil artefak 1 dan 2, lalu artefak 3. Artefak lenyap dan skor
  bertambah, tetapi bunyinya tidak terdengar & Sedang & Selesai \\
2 & Mekanik & Melompat di tepi rak, lalu tekan Space lagi sebelum mendarat.
  Pemain melompat dua kali & Sedang & Selesai \\
3 & Keseluruhan & Berjalan ke sudut kanan bawah rak kanan dan melompat ke
  dinding. Pemain tersangkut di celah dangkal & Berat & Selesai \\
4 & Tampilan & Mendarat, lalu segera berjalan. Animasi diam bertahan sekitar
  satu detik & Ringan & Terbuka \\
5 & Tampilan & Mainkan pada jendela 4:3. Teks nyawa bertumpuk dengan ikon
  ketiga & Ringan & Ditunda \\
\bottomrule
\end{tabularx}
}

\subsection*{Tabel susunan demo}

{\footnotesize
\renewcommand{\arraystretch}{1.2}
\begin{tabularx}{\linewidth}{@{}C{1.6cm} L{3.0cm} Y@{}}
\toprule
\textbf{Waktu} & \textbf{Bagian} & \textbf{Isi} \\
\midrule
1 menit & Pembukaan & Nama, jenis, dan satu kalimat inti permainan \\
4--6 menit & Demo langsung & Permainan dijalankan dari awal sampai tujuan \\
1--2 menit & Temuan pengujian & Satu atau dua temuan terpenting dan
  perbaikannya \\
1 menit & Penutup & Kekurangan yang diketahui dan rencana berikutnya \\
\bottomrule
\end{tabularx}
}


\checkpoint{05}{\raggedright Folder checkpoint memuat seluruh skrip dari
\file{Scripts/}, \file{scene.md}, \file{gdd.md},
\file{checklist-fungsional.md}, \file{lembar-playtest.md},
\file{log-bug.md}, dan \file{naskah-demo.md}. Build dan berkas biner tidak
disertakan; hasilkan kembali dari project sesuai \file{scene.md}.}

% ---------------------------------------------------------
\bagian{Kesalahan Umum}

\begin{galatbox}
\begin{itemize}[leftmargin=1.4em, itemsep=0.35em]
  \item \textbf{Build hanya berjalan di komputer pembuatnya.}
        Gejala: berkas \file{.exe} dibuka di komputer lain lalu gagal atau
        layar kosong. Periksa apakah scene sudah masuk \menu{Add Open Scenes},
        apakah seluruh isi folder keluaran disalin (termasuk folder
        \file{Data}), dan apakah platform yang dipilih sesuai.
  \item \textbf{Checklist terlalu umum atau terlalu panjang.}
        Gejala: butir seperti ``permainan terasa baik'' tidak dapat dijawab
        lulus atau tidak lulus, atau daftar 60 butir tidak pernah selesai.
        Periksa setiap butir terhadap GDD; ubah menjadi pernyataan yang dapat
        diperiksa, dan batasi sekitar 20--30 butir.
  \item \textbf{Playtest menghasilkan catatan yang menyenangkan tetapi kosong.}
        Gejala: penguji menulis ``cukup jelas'' dan Anda tidak memperoleh
        temuan. Periksa cara Anda memandu: berhenti memberi petunjuk, catat
        waktu dan tempat berhenti, dan ajukan pertanyaan sesudah sesi, bukan
        selama permainan.
  \item \textbf{Catatan bug tidak dapat diulang.}
        Gejala: ``bunyi kadang hilang'' tanpa urutan tindakan. Periksa apakah
        catatan memuat apa yang terjadi, apa yang seharusnya terjadi, dan
        langkah yang menyebabkannya. Sempitkan dugaan dengan mencoba satu
        tindakan setiap kali.
  \item \textbf{Janji GDD dihapus agar laporan tampak sesuai.}
        Gejala: GDD berubah tanpa catatan alasan, dan fitur yang tidak
        dikerjakan hilang dari dokumen. Periksa kembali versi GDD lama;
        cantumkan keadaan sesungguhnya dan alasan perubahannya.
  \item \textbf{Demo melebihi waktu atau gagal secara teknis.}
        Gejala: presentasi berhenti di tengah karena build tidak berjalan, atau
        waktu habis sebelum temuan disampaikan. Periksa latihan berwaktu,
        video cadangan, dan build cadangan pada dua tempat.
\end{itemize}
\end{galatbox}

% ---------------------------------------------------------
\bagian{Lembar Kerja}

Isi tabel berikut, lalu simpan sebagai bagian dari catatan praktikum Anda.

{\small
\renewcommand{\arraystretch}{1.4}
\begin{tabularx}{\linewidth}{@{}C{0.8cm} L{6.2cm} Y@{}}
\toprule
\textbf{No} & \textbf{Hasil yang dicatat} & \textbf{Keterangan Anda} \\
\midrule
1 & Jumlah butir checklist dan sebarannya pada tiga kelompok & \\[0.5cm]
2 & Jumlah butir yang lulus pada putaran 1 dan putaran 2 & \\[0.5cm]
3 & Temuan berat beserta perbaikannya & \\[0.9cm]
4 & Nama dan pengalaman dua penguji playtest & \\[0.5cm]
5 & Waktu penyelesaian dan titik berhenti setiap penguji & \\[0.9cm]
6 & Perbedaan GDD dan build yang Anda temukan & \\[0.9cm]
7 & Durasi demo hasil latihan terakhir & \\[0.5cm]
\bottomrule
\end{tabularx}
}

% ---------------------------------------------------------
\bagian{Checklist Sesi}

\begin{checklistbox}
\begin{checklist}
  \item Saya dapat menunjukkan build yang berjalan pada komputer lain tanpa
        Unity, beserta suara dan kontrolnya.
  \item Saya dapat menunjukkan GDD satu halaman dan menjelaskan setiap bagiannya.
  \item Saya dapat menunjukkan checklist fungsional yang butirnya diturunkan
        dari GDD dan terbagi menjadi tiga kelompok.
  \item Saya dapat menunjukkan lembar playtest yang terisi oleh dua penguji,
        lengkap dengan waktu dan titik berhenti.
  \item Saya dapat menunjukkan log bug yang setiap catatannya dapat diulang
        dan memiliki berat serta status.
  \item Saya dapat menyajikan demo lima sampai sepuluh menit dan menjelaskan
        satu temuan beserta perbaikannya.
\end{checklist}
\end{checklistbox}

% ---------------------------------------------------------
\bagian{Refleksi: Dari Python ke Unity}

Sepanjang Blok I dan Blok II Anda menulis sendiri siklus permainan,
\code{delta time}, gravitasi, tabrakan, animasi sederhana, pencarian jalur A*,
dan pengaturan bunyi menurut jarak. Pada Unity, pekerjaan yang sama diserahkan
kepada mesin: \code{Update} dan \code{FixedUpdate} dipanggil oleh
\istilah{game engine}, \istilah{Rigidbody2D} mengurus fisika dan tabrakan,
\istilah{Animator} memilih animasi, \istilah{AudioSource} mengurus pemutaran
dan pelemahan bunyi menurut jarak, dan \istilah{Canvas} mengurus tata letak
antarmuka. Yang tetap menjadi pekerjaan Anda adalah hal yang tidak dapat
diwakilkan: merancang aturan, menyusun keadaan, menghubungkan peristiwa,
memilih nilai, dan menguji hasilnya. Perbedaan itu bukan penilaian mana yang
lebih baik, melainkan keterangan tentang pekerjaan mana yang perlu Anda kuasai
sendiri dan pekerjaan mana yang sebaiknya Anda serahkan kepada perkakas.

% ---------------------------------------------------------
\bagian{Penutup}

Pada minggu keenam belas, UAS menilai proyek yang dihasilkan modul ini.
Yang diserahkan adalah \istilah{build} yang dapat dijalankan,
\istilah{project source}, GDD, checklist fungsional, lembar \istilah{playtest},
\istilah{log bug}, naskah demo, dan laporan perbaikan. Uji fungsional oleh
dosen memakai checklist yang serupa dengan yang Anda susun, sehingga butir
yang jujur dan dapat diperiksa akan menolong Anda sendiri.

Bawalah build pada dua tempat, salinan project source, dan seluruh dokumen
dalam bentuk cetak maupun berkas. Periksa sekali lagi bahwa build berjalan
pada komputer lain, bahwa tidak ada temuan berat yang masih terbuka, dan bahwa
demo Anda selesai dalam waktu yang ditentukan. Kekurangan yang Anda ketahui
dan sebutkan dengan jujur lebih bernilai daripada permainan yang tampak tanpa
cacat tetapi tidak dapat dipertanggungjawabkan.
